\documentclass[../../main.tex]{subfiles}
\graphicspath{{../../image/}{../../figure/}}
\begin{document}

\section{模型平均}\label{section:ch09-模型平均}

如果我们训练了几个不同模型来解决相同问题，而不是尝试选择单一最佳模型，那么通常可以通过平均单个模型所做预测来提高泛化能力。这种模型整合方式有时称为\textbf{委员会}或\textbf{集成}（ensemble）。对于产生概率输出的模型，预测分布是每个模型预测的平均值：
\begin{align}
p(t\mid \mathbf{x}) = \frac{1}{L}\sum_{l=1}^{L} p_l(t\mid \mathbf{x}),
\label{align:ch09-ensemble}
\end{align}
其中 $p_l(t\mid \mathbf{x})$ 是模型 $l$ 的输出，$L$ 是模型总数。

这种平均过程的动机源于偏差和方差之间权衡。回想图~4.7：当使用正弦数据训练多个多项式，然后对得到的函数进行平均时，方差项产生的贡献几乎会被抵消，从而提升预测性能（参见~4.3~节）。当然在实践中只有一个数据集，因此必须找到一种方法，以便在集成模型（一些旨在解决相同问题的模型集合）内部不同模型之间引入可变性。

可以使用\textbf{自助}（bootstrap）数据集，其中多个数据集创建过程如下。假设原始数据集由 $N$ 个数据点 $\mathbf{X}=\{\mathbf{x}_1,\cdots,\mathbf{x}_N\}$ 组成。可以通过从 $\mathbf{X}$ 中有放回地抽样 $N$ 个数据点来创建新数据集 $\mathbf{X}_B$，这样 $\mathbf{X}$ 中一些数据点可能在 $\mathbf{X}_B$ 中出现多次，而 $\mathbf{X}$ 中其他数据点则可能不包含在 $\mathbf{X}_B$ 中。这个过程可以重复 $L$ 次，从而生成 $L$ 个大小为 $N$ 的数据集，其中每个数据集都是通过从原始数据集 $\mathbf{X}$ 中采样得到。每个数据集可以用来训练一个模型，并对得到的模型预测进行平均。这个过程称为\textbf{自助聚合}（bootstrap aggregation）或\textbf{袋装}（bagging）~\citep{Breiman1996}。

形成集成模型的另一种方法是使用原始数据集训练多个具有不同架构的模型。

我们可以通过考虑一个仅涉及输入向量 $\mathbf{x}$ 和单输出变量 $y$ 的回归问题来分析集成预测的好处。假设我们有一组训练好的模型 $y_1(\mathbf{x}),\cdots,y_L(\mathbf{x})$，构造由下式给出的集成模型预测：
\begin{align}
y_{\mathrm{ens}}(\mathbf{x}) = \frac{1}{L}\sum_{l=1}^{L} y_l(\mathbf{x}).
\label{align:ch09-ens-pred}
\end{align}

如果试图预测的真实函数由 $h(\mathbf{x})$ 给出，那么每个模型输出都可以写成真实值加上一个误差的形式：
\begin{align}
y_l(\mathbf{x}) = h(\mathbf{x})+\epsilon_l(\mathbf{x}).
\label{align:ch09-ens-error}
\end{align}

平均平方和误差则采用以下形式：
\begin{align}
\mathbb{E}_{\mathbf{x}}\!\left[\{y_l(\mathbf{x})-h(\mathbf{x})\}^2\right] = \mathbb{E}_{\mathbf{x}}[\epsilon_l(\mathbf{x})^2],
\label{align:ch09-ens-mse}
\end{align}
其中 $\mathbb{E}_{\mathbf{x}}[\cdot]$ 表示关于输入向量 $\mathbf{x}$ 分布的期望。因此集合中模型分别预测产生的平均误差为
\begin{align}
\mathbb{E}_{\mathrm{avg}} = \frac{1}{L}\sum_{l=1}^{L}\mathbb{E}_{\mathbf{x}}[\epsilon_l(\mathbf{x})^2].
\label{align:ch09-ens-avg}
\end{align}

类似地，来自集成模型~\eqref{align:ch09-ens-pred}~的期望误差由下式给出：
\begin{align}
\mathbb{E}_{\mathrm{ens}} = \mathbb{E}_{\mathbf{x}}\!\left[\left(\frac{1}{L}\sum_{l=1}^{L}\epsilon_l(\mathbf{x})\right)^2\right].
\label{align:ch09-ens-expected}
\end{align}

如果假设误差均值为零且不相关，则有
\begin{align}
\mathbb{E}_{\mathbf{x}}[\epsilon_l(\mathbf{x})] = 0,
\label{align:ch09-ens-meanzero}
\end{align}
\begin{align}
\mathbb{E}_{\mathbf{x}}[\epsilon_m(\mathbf{x})\epsilon_l(\mathbf{x})] = 0, \quad m\neq l.
\label{align:ch09-ens-uncorr}
\end{align}

于是可以得到（见习题~9.14）：
\begin{align}
\mathbb{E}_{\mathrm{ens}} = \frac{1}{L}\mathbb{E}_{\mathrm{avg}}.
\label{align:ch09-ens-result}
\end{align}

这个结果表明，只需要对 $L$ 个不同版本模型进行平均，就可以将模型平均误差降低到原来 $1/L$。遗憾的是它依赖于一个关键假设，即各个模型造成的误差是不相关的。在实践中误差通常高度相关，总体误差减小通常要小得多。然而可以证明，预期的集成模型误差不会超过集合中模型的预期误差，因此
\begin{align}
\mathbb{E}_{\mathrm{ens}} \le \mathbb{E}_{\mathrm{avg}}.
\label{align:ch09-ens-bound}
\end{align}

另一种称为\textbf{提升}（boosting）~\citep{FreundSchapire1996}的模型组合方法，将多个"基"分类器组合起来，产生一种性能明显优于任何基分类器的集成模型。即使基分类器表现只比随机好一点点，提升也可以给出好结果。提升方法和前面讨论的集成模型方法（如袋装）的主要区别在于：基分类器是按顺序训练的，且每个基分类器都是使用数据集的加权形式进行训练的，其中与每个数据点相关的权重系数取决于前一个分类器的性能。特别是被其中一个基分类器误分类的点，在训练序列中下一个分类器时会被赋予更大权重。一旦所有分类器都训练完毕，就可以对它们预测通过加权多数投票方式进行整合（见习题~9.15）。

在实践中所有模型组合方法的主要缺点是必须训练多个模型，然后必须对所有模型预测进行评估，因而增加了训练和推理计算开销，开销大小取决于具体应用场景。

\subsection{Dropout}\label{section:ch09-dropout}

\textbf{Dropout}~\citep{Srivastava2014}是一种使用广泛且非常有效的正则化形式，它可以视为一种在不需要单独训练多个模型的情况下，对指数级数量规模模型隐式地进行模型平均的方法。它具有广泛适用性，并且计算开销低。

dropout 核心思想是在训练过程中随机删除网络中的节点，包括它们之间的连接。每当一个新数据点进入网络，模型就会随机选择一些节点并忽略它们。dropout 可以同时应用于隐藏节点和输入节点，但不能应用于输出节点，这相当于将被删节点输出设置为零。它可以通过定义一个掩码向量 $R_i\in\{0,1\}$ 来实现：对于数据点 $n$，该掩码将被乘以非输出节点 $i$，其值以概率 $p$ 设置为 $1$。对于隐藏节点 $p=0.5$ 很有效；而对于输入节点，通常使 $p=0.8$。

在训练过程中，当每个数据点进入网络时就会创建一个新掩码，并在该剪枝网络上应用前向和后向传播步骤来创建误差函数梯度，然后用于更新权重，例如采用随机梯度下降法。如果数据点被分组成多个小批次，则对梯度在每个小批次（参见~7.2.4~小节）中的数据点上进行平均后，再应用权重更新。对于具有 $M$ 个非输出节点的网络，有 $2^M$ 个剪枝网络，因此在训练期间只会考虑这些网络中的一少部分。这与传统集成方法不同：在传统集成方法中集成中的每个网络都是独立训练至收敛的。它们之间另一个区别是，指数级数量规模的网络隐式地用 dropout 进行训练，它们不是独立的，而是与完整网络共享它们的参数值，因此剪枝网络间参数也是彼此共享的。

需要注意的是，使用 dropout 训练后，单个参数更新包含较大噪声，因此训练可能需要花费更长时间。此外由于误差函数本身是有噪声的，仅通过在训练过程中减小误差函数很难确保优化算法正常工作。

一旦训练完成，原则上就可以通过应用集成规则~\eqref{align:ch09-ensemble}~来进行预测，这种情况下的集成规则如下：
\begin{align}
p(t\mid \mathbf{x}) = \sum_R p(R)\,p(t\mid \mathbf{x},R),
\label{align:ch09-dropout-ensemble}
\end{align}
其中求和是在掩码的指数级空间上进行的，$p(t\mid \mathbf{x},R)$ 是来自具有掩码 $R$ 的网络的预测分布。因为这种求和难以处理，所以可以通过对少量掩码进行采样来近似。在实践中只需要~10~或~20~个掩码就足以获得良好结果。这一过程称为蒙特卡洛 dropout。

一种更简单方法是在不掩盖训练好的网络节点情况下进行预测，并重新缩放网络中权重，使每个节点在测试期间的预期输入与训练期间输入大致相同，以补偿在训练中会忽略一部分节点的事实。因此如果一个节点在训练期间以概率 $p$ 的形式存在，那么在测试期间在使用网络进行预测之前，需要对来自该节点的输出权重乘以 $p$。

应用 dropout 的另一种动机来自贝叶斯学派。在按照贝叶斯学派观点处理问题时，我们将通过对所有 $2^M$ 个网络模型进行平均来预测，其中每个网络模型都由其后验概率加权（参见~2.6~节）。这样做计算成本很高，无论是在评估后验概率的训练期间，还是在计算加权预测的测试期间。dropout 通过给每个可能模型赋予相同权重来近似这个模型平均。

dropout 具有良好表现背后的更深层次原因，在于 dropout 在减少过拟合方面发挥了重要作用。在一个标准神经网络中，参数可以调整为错误适应单个数据点，其隐藏节点被过度特化：给定其他节点输出，每个节点会调整其权重以最小化误差，导致节点以一种可能无法泛化到新数据的方式共同调整参数。使用 dropout，每个节点不能依赖其他特定节点存在，而必须在广泛场景下做出有用贡献，从而减少共同调整和特化。对于使用最小二乘法训练的简单线性回归模型，dropout 正则化等价于二次正则化的一种修改形式（见习题~9.18）。

\end{document}
